% ===========================================================================
%  user/04-streams/00-chapter.tex
% ===========================================================================
\chapter{Streams}\label{cap:streams}

A \voce{stream} is the result of a strategy: the series of its gains and
losses over time. The \voce{Streams} page shows them all, one card each.

\figuraSchermata{1.0}{figure/schermate/streams.png}
  {The list of results. Each card carries the performance, the Sharpe ratio,
   the return and the maximum loss.}
  {streams}

\section{Running a backtest}\label{sez:run}

The \voce{Run backtest} tab, next to \voce{Streams}, lists everything that
can be run: the installed strategies and approaches, each with its own
description and a button.

\figuraSchermata{1.0}{figure/schermate/run.png}
  {What can be run. \pulsante{Run All} queues them all, one after the
   other.}
  {run}

\begin{pratica}{Running it, and then}
\begin{enumerate}
  \item Press the button of the strategy you want.
  \item The page shows the log live: you see what it is doing, line by line.
  \item You can close the page. The job carries on, and the bell at the top
        right tells you when it is done.
  \item The result appears in \voce{Streams}.
\end{enumerate}
\end{pratica}

\begin{attenzione}{One job at a time}
Backtests run in a queue, not in parallel, even if you press several
buttons. It is not a limit to work around: two calculations at the same time
would fight over the same files and produce mixed-up results.
\end{attenzione}

\section{The page of a result}\label{sez:stream-detail}

Clicking on a stream opens its page.

\figuraSchermata{1.0}{figure/schermate/stream.png}
  {The page of a stream. At the top the robustness verdict, below it the
   metrics, then the curve.}
  {stream}

\subsection{Robustness: the panel to read first}

It is the most important part of the page, and it is the one that answers
the question \enquote{did this strategy work, or was it lucky?}. The score on
the right sums up the five numbers in one.

\begin{center}
\begin{tabular}{@{}L{42mm}L{96mm}@{}}
\toprule
\textbf{Indicator} & \textbf{What it tells you} \\
\midrule
\rigaTenue
\voce{P(Sharpe > 0)} & The probability that the result is not noise. Below
  \perc{95} the strategy has proved nothing. \\
\voce{Monte Carlo p-value} & The same question asked another way: shuffling
  the data, how often would a lucky case have done just as well. The smaller,
  the better. \\
\rigaTenue
\voce{Min. track record} & How long this strategy would need to run for its
  result to be credible. A value longer than the period tested means: there
  is not enough history yet. \\
\voce{OOS Sharpe instability} & How much the Sharpe ratio changes from one
  period to the next. High means the strategy worked in some periods and not
  in others. \\
\rigaTenue
\voce{Losing OOS periods} & The percentage of periods closed at a loss. \\
\bottomrule
\end{tabular}
\end{center}

\begin{perche}{Why a \enquote{Doubtful} verdict is not an error}
It means the data are not enough to tell this strategy apart from chance.
Not that it is wrong: that it has not been proved. It is the honest answer
over a short period, and the way to change it is to have more history ---
not to change the parameters until the verdict improves, which is the most
common way of fooling yourself.
\end{perche}

\subsection{The metrics}

Ten numbers, unless you choose otherwise (see below). Besides those already
seen on the home screen (\vediCap{overview}):

\begin{description}[leftmargin = 1.4em, style = nextline, itemsep = 4pt]
  \item[\voce{Beta vs SPY}] How closely the strategy follows the market.
    Close to zero: it does something of its own. Close to one: it moves like
    the index.
  \item[\voce{Sortino}] Like the Sharpe ratio, but it counts only the
    downward swings. Rising fast is not a flaw.
  \item[\voce{Calmar}] Annual return divided by maximum loss. It answers
    \enquote{how much does it earn for each point of pain}.
  \item[\voce{Hit rate}] The percentage of periods in profit.
  \item[\voce{N obs}] How many observations. Few observations make
    everything else less reliable.
\end{description}

\subsection{Choosing the metrics}\label{sez:stream-metrics}

Every metric on the page can be hidden: the numbers, the robustness panel,
the charts, the tables, and the metrics you wrote yourself in Studio. Move
the pointer over one and a small eye appears in its corner: pressing it
hides that metric, and the message that follows has an \pulsantesec{Undo}.
To see the whole list, press \pulsantesec{Choose metrics}, at the top of
the metrics.

\figuraSchermata{1.0}{figure/schermate/stream-metrics.png}
  {\pulsantesec{Choose metrics}. Numbers, charts and tables, the default
   ones and yours, each with its tick. Next to the button, how many are
   shown.}
  {stream-metrics}

\begin{itemize}
  \item The choice holds for \emph{every} stream: you make it once.
  \item \pulsantesec{Show all}, \pulsantesec{Hide all} and
    \pulsantesec{Default} change everything at once. \voce{Default} is how
    the page looks before any choice: the default metrics and your active
    ones.
  \item A metric you create after making your choice appears by itself; you
    hide it like the others if you do not want it.
  \item Your metrics are computed only when they are shown, on the results
    of the stream and with its benchmark.
\end{itemize}

\noindent
One of your metrics can show a word instead of a value. \voce{Data not
available}: it reads data that is not in the system yet --- load it in
\percorso{Settings \freccia Data} and the metric works. \voce{Error in the
code}: the metric has to be fixed in Studio, and the word takes you there.
\voce{No results yet}: the stream has never been run.

\begin{pratica}{Adding a metric from here}
Move the pointer after the numbers, or at the very bottom of the page, below
the charts and the tables: \pulsantesec{New number metric} and
\pulsantesec{New chart metric} appear. They open the \voce{Metrics} section
of Studio with the output already chosen, a template and a free name; save,
and the metric is on the stream page (\vediSez{studio-metrics}) --- your
numbers after the default ones, your charts at the bottom, just above
\pulsantesec{New chart metric}. The pencil on one of your metrics opens it in
Studio, and Studio has a link back to the stream.
\end{pratica}

\subsection{Arranging charts and tables}\label{sez:stream-layout}

Charts and tables --- the default ones and yours --- are panels you arrange
like windows on a desktop. They stay in rows, up to three side by side.

\figuraSchermata{1.0}{figure/schermate/stream-layout.png}
  {Moving a panel. Where it was, a dashed outline; under the pointer, a small
   copy; the blue frame shows where it will land --- here the right third of
   the row, next to two other panels.}
  {stream-layout}

\begin{itemize}
  \item \textbf{Move.} Take a panel by its title bar --- the strip with its
    name --- and drag it. A blue frame shows where it will go: next to the
    panels of a row (\voce{Left half}, \voce{Middle third}\ldots), or a
    line \voce{New row} above or below a row. Let go, and the panels slide
    into place. \voce{Esc} cancels. Near the top or the bottom of the window
    the page scrolls by itself.
  \item \textbf{Make it bigger.} Pull the bottom edge of a panel down: the
    whole row grows, because panels side by side keep the same height.
    Between two panels side by side, the divider changes their widths.
    Double-click the edge or the divider: back to the normal size.
  \item \textbf{Maximize.} Double-click a panel's title: it takes the whole
    width and the height of the window. Double-click again: it goes back to
    where it was, as it was.
  \item \textbf{With the keyboard.} Among the buttons that appear on a
    panel there is one with four arrows. With the focus on it, the arrow keys
    move the panel, \voce{Shift} with the arrows change its height and width,
    \voce{Enter} maximizes it and brings it back.
\end{itemize}

\noindent
Like the choice of the metrics, the arrangement holds for every stream.
\pulsantesec{Reset layout}, in \pulsantesec{Choose metrics}, brings back
the original order: one panel per row, at the normal size. One of your
metrics that draws several charts makes one panel for each; a new chart of
yours appears at the bottom, just above \pulsantesec{New chart metric}, and
you move it where you want. In a narrow window the panels go one below the
other.

\subsection{The three buttons at the top}

\begin{itemize}
  \item \pulsante{Open the code in Studio} --- opens the artifact in the
    editor.
  \item \pulsantesec{Run backtest} --- runs the calculation again.
  \item \pulsantesec{Use in Ensemble} --- takes this stream into the project
    builder (\vediCap{projects}).
\end{itemize}

\section{View: seeing what it did}\label{sez:view}

The \voce{View} item in the menu shows prices with the strategy's trades on
top: where it went in, where it came out, and the indicators it was
watching.

\figuraSchermata{1.0}{figure/schermate/visual.png}
  {The chart view. You choose the strategy and the instrument, and overlay
   the indicators.}
  {visual}

\noindent
It is for something the numbers do not say: \emph{whether the strategy is
doing what you thought}. A curve that rises nicely can hide twenty sensible
trades and a single huge lucky one.
